\documentclass[a4paper,10pt]{article}
\usepackage[margin=1.5cm]{geometry}
\usepackage{multicol}
\usepackage{hyperref}
\usepackage{enumitem}
\setlength{\parindent}{0pt}

\begin{document}

\begin{center}
    {\Large \textbf{Formulario de Comandos Linux}}\\
\end{center}


\section*{Hardware Information}
\begin{center}
    \begin{tabular}{| c | c |}
        \hline
        Comando & Función \\
        \hline\hline
        \texttt{lscpu} & Muestra información de la CPU \\
        \hline
        \texttt{lsblk} & Muestra información sobre \textbf{block devices} (HDD O SSD) \\
        \hline
        \texttt{lspci -tv} & Muestra los dispositivos \textbf{PCI} (Peripheral Component Interconnec)\\ & en forma de arbol (-tv) \\
        \hline
        \texttt{lsusb -tv} & Despliega los dispositivos USB en formato de arbol (-tv) \\
        \hline
        \texttt{lshw} & Muestra la configuración del Hardware \\
        \hline
        \texttt{cat /proc/cpuinfo} & Muestra información detallada de la CPU. \\
        \hline
        \texttt{cat /proc/meminfo} & Información detallada de la memoria del sistema. \\
        \hline
        \texttt{cat /proc/mounts} & Muestra información de \textbf{mounted files}. \\
        \hline
        \texttt{free -h} & Muestra memoria ocupada y disponible en formato humano (-h). \\
        \hline
        \texttt{sudo dmidecode} & Muestra información del HW desde la \textbf{BIOS}. \\
        \hline
        \texttt{sudo hdparm -i /dev/[device name]} & Despliega información (-i) del disco. \\
        \hline
        \texttt{sudo hdparm -tT /dev/[device name]} & Hace una prueba de velocidad de lectura (-tT) al disco. \\
        \hline
        \texttt{sudo badblocks -s /dev/[device name]} & Prueba para encontrar bloques no leibles en el disco. \\
        \hline
        \texttt{sudo fsck /dev/[device name]} & Hace un chequéo al disco. \\
        \hline
    \end{tabular}
\end{center}


\section*{Comandos de busqueda (Searching commands)}
\begin{center}
    \begin{tabular}{| c | c |}
        \hline
        Comando & Función \\
        \hline\hline
        \textbf{find [path] $-$name [search $\_$ pattern]} & Encuentra archivos y directorios que cumplan con\\ & el patrón dentro del directorio. \\
        \hline
        \textbf{find [path] $-$size [+100M]} & Muestra archivos y directorios con tamaño mayor o igual al dado. \\
        \hline
        \textbf{grep [search$\_$pattern] [file$\_$name]} & Busca un patrón dentro de un archivo. \\
        \hline
        \textbf{grep -r [search$\_$pattern] [directory$\_$name]} & Busca un patrón dentro de archivos de un direcotio recursivamente \\
        \hline
        \textbf{locate [name]} & Localiza TODOS los archivos y directorios relacionados a un nombre. \\
        \hline
        \textbf{which [command]} & Busca la ruta del comando dentro de la variable de entorno $\%$ PATH \\
        \hline
        \textbf{whereis [command]} & Encuentra la fuente, el binario y el manual de un comando \\
        \hline
        \textbf{(g)awk '[search$\_$pattern] {print \$0}' [file$\_$name]} & Imprime todas las lineas que contengan al patrón. \\
        \hline
        \textbf{sed 's/[old$\_$text]/[new$\_$text]/' [file$\_$name]} & Encuentra y reemplaza texto en un archivo. \\
        \hline
    \end{tabular}
\end{center}


\section*{Archivos y Directorios}
\subsection*{ls}
\begin{itemize}[leftmargin=*]
    \item \texttt{ls} — listar archivos
    \item \texttt{cd dir} — cambiar directorio
    \item \texttt{pwd} — ruta actual
    \item \texttt{mkdir dir} — crear directorio
    \item \texttt{rm -r dir} — borrar directorio
    \item \texttt{cp src dst} — copiar
    \item \texttt{mv src dst} — mover/renombrar
\end{itemize}

\section*{Permisos}
\begin{itemize}[leftmargin=*]
    \item \texttt{chmod 755 file}
    \item \texttt{chown user:group file}
    \item \texttt{ls -l}
\end{itemize}

\section*{Procesos}
\begin{itemize}[leftmargin=*]
    \item \texttt{ps aux}
    \item \texttt{top / htop}
    \item \texttt{kill PID}
    \item \texttt{kill -9 PID}
\end{itemize}

\section*{Red}
\begin{itemize}[leftmargin=*]
    \item \texttt{ip a}
    \item \texttt{ping host}
    \item \texttt{ss -tulpn}
    \item \texttt{curl url}
\end{itemize}

\section*{Compresión}
\begin{itemize}[leftmargin=*]
    \item \texttt{tar -czf file.tar.gz dir}
    \item \texttt{tar -xzf file.tar.gz}
    \item \texttt{zip -r file.zip dir}
\end{itemize}


\section*{Fuentes}
\begin{itemize}
    \item https://phoenixnap.com/kb/linux-commands-cheat-sheet
\end{itemize}

\end{document}
